\section{Introduction}
\label{sec:intro}

Can a machine simulate the full spectrum of human possibility? As Large Language Models (LLMs) transition from passive text generators to active "generative agents" \cite{park2023generative}, they are being tasked with representing human behavior in increasingly complex social, economic, and policy-driven simulations \cite{horton2023homo}. In these roles, LLMs do not just predict tokens; they inhabit characters, make decisions, and interact within simulated worlds. However, despite their widespread use, we still lack a formal "Theory of Simulation" that explains how these models navigate the space of potential human actions.

Understanding LLMs as simulators is crucial because it determines our ability to use them for risk assessment and "what-if" analysis. If an LLM merely replicates the most frequent or "average" human behavior, it becomes a poor tool for identifying the "black swan" events or extreme tail outcomes that often drive systemic change. A valid simulator of human behavior must be able to capture the underlying psychological and cultural factors that lead to diverse outcomes, ranging from submissive cooperation to aggressive confrontation.

{\bf What is missing in existing work?} Most current research focuses on the *accuracy* or *believability* of LLM simulations---measuring how closely they match the single most likely human response in a given dataset. This "point-estimate" view ignores the fact that human behavior is inherently multiversal; the same person in the same situation might choose different paths based on subtle internal or external shifts. There is a significant gap in exploring the **full distribution** of simulated outcomes. Specifically, can LLMs simulate the "tails" of human behavior? Does their training on the vast superset of human text allow them to simulate a "superset" of outcomes that exceeds any single individual's behavior but remains within the bounds of collective human possibility?

In this paper, we propose a manifold-based theory of simulation. We conceptualize the LLM as a high-dimensional manifold representing the total sum of documented human behavior. We hypothesize that personas and scenarios act as constraints that project the probability mass onto specific sub-regions of this manifold. To test this, we introduce the \method framework, where we roll out multiple parallel continuations from a single starting state to map the distribution of potential futures (see \figref{fig:method}).

Our results show that \gpt generates a wide distribution of plausible outcomes, with entropy levels varying significantly based on persona conditioning. We find a striking "alignment bias" in generic personas, where 98\% of outcomes collapse into a single "safe" mode of behavior. However, by using strategic persona conditioning and explicit tail-prompting, we can recover the model's ability to traverse the full superset of human behavior, including rare but highly plausible outcomes like a "humorous Oscar joke" in a high-tension social dilemma.

We summarize our primary contributions as follows:
\begin{itemize}[leftmargin=*,itemsep=0pt,topsep=0pt]
    \item We propose a **Manifold-based Theory of Simulation** for LLMs, moving from point-estimate evaluations to multiversal distributional analysis.
    \item We introduce the **\method framework** and demonstrate its utility in mapping the "outcome manifold" of social conflict across 200+ parallel simulations.
    \item We empirically characterize **alignment-induced mode collapse** in generic personas, where behavioral entropy drops by 86\% compared to strategic personas.
    \item We demonstrate **cultural and factor grounding**, showing that LLMs accurately shift their behavior manifolds to accommodate cultural norms such as Japanese "wa" (harmony) vs. NYC "disruptive honesty."
\end{itemize}
